% sections/03_engineering.tex
% 工程层分析：Q4--Q6

\section{工程层：发射、供电与可靠性}

科学层确认了散热在物理极限内的可行性，但工程层需要回答：在当前技术水平下，将这些物理方案付诸实施的可靠性如何？本章聚焦三个关键工程维度：发射成本的历史趋势、轨道光伏供电链路的实际效率、以及GPU在轨可靠性的数据缺口。

\subsection{发射成本：Wright学习率的现实检验}

SpaceX的Falcon 9将LEO发射成本从航天飞机时代的\$65,000/kg降至约\$2,720/kg（70\%降幅），但此后下降速度显著放缓。表~\ref{tab:launch_cost} 整理了关键里程碑。

\begin{table}[H]
\centering
\footnotesize
\caption{LEO发射成本里程碑}
\label{tab:launch_cost}
\begin{tabular}{lcc}
\toprule
\textbf{运载器/阶段} & \textbf{LEO成本（\$/kg）} & \textbf{来源} \\
\midrule
航天飞机（1981--2011） & $\sim$65,000 & NASA\cite{nasalaunch2018} \\
Atlas V（2002--） & $\sim$13,000 & NASA \\
Falcon 9 v1.0（2010） & $\sim$4,300 & SpaceX \\
Falcon 9 FT（2015--） & $\sim$2,720 & SpaceX \\
Falcon Heavy（2018--） & $\sim$1,400 & SpaceX \\
Starship（理论预测） & $\sim$200（目标） & SpaceX/FAA\cite{faa2025boca} \\
\bottomrule
\end{tabular}
\end{table}

图~\ref{fig:launch_cost} 以半对数坐标直观呈现发射成本40年间下降272倍的轨迹——以及Starship目标与物理地板之间的剩余空间。

\begin{figure}[H]
\centering
% launch_cost_curve.tex — 发射成本半对数历史曲线
\begin{tikzpicture}
\begin{semilogyaxis}[
    width=0.95\textwidth,
    height=7cm,
    xlabel={年份},
    ylabel={\$/kg to LEO（对数坐标）},
    grid=both,
    grid style={draw=gray!20},
    minor grid style={draw=gray!10},
    xtick={1980,1990,2000,2002,2010,2015,2018,2020,2023,2026,2030},
    xticklabels={1980,1990,2000,2002,2010,2015,2018,2020,2023,2026,2030},
    ytick={10,30,100,500,2000,5000,20000,70000},
    yticklabels={10,30,100,500,2{,}000,5{,}000,20{,}000,70{,}000},
    ymin=8, ymax=100000,
    xmin=1978, xmax=2032,
    legend style={
        at={(0.5,-0.22)},
        anchor=north,
        legend columns=3,
        draw=none,
        font=\footnotesize
    },
    mark options={solid},
    line width=1.2pt,
    every node near coord/.append style={font=\footnotesize}
]

% 历史数据点
\addplot[color=primary, mark=*, mark size=2.5pt]
    coordinates {
        (1981,54500)(1985,54500)(1991,54500)(2001,54500)(2011,54500)
    };
\addlegendentry{航天飞机 (1981--2011)}

\addplot[color=primary, mark=square*, mark size=2.5pt]
    coordinates {
        (2002,13000)(2005,13000)(2010,13000)(2015,13000)
    };
\addlegendentry{Atlas V (2002--)}

\addplot[color=primary, mark=triangle*, mark size=2.5pt]
    coordinates {
        (2010,4300)(2015,2720)(2018,2400)(2020,2400)(2023,2400)
    };
\addlegendentry{Falcon 9 (2010--)}

\addplot[color=primary, mark=diamond*, mark size=2.5pt]
    coordinates {
        (2018,1400)(2020,1400)(2023,1400)
    };
\addlegendentry{Falcon Heavy (2018--)}

% Starship 目标区间
\addplot[color=accent, dashed, line width=1.0pt, mark=none]
    coordinates {
        (2024,500)(2026,200)
    };
\addlegendentry{Starship 目标}

\addplot[color=accent, opacity=0.2, mark=none]
    coordinates {
        (2026,100)(2030,50)
    } \closedcycle;
\addlegendentry{Starship 乐观区间}

% 物理地板
\addplot[color=highlight, dashed, line width=1.0pt, mark=none]
    coordinates {
        (1978,25)(2032,25)
    };
\node[highlight, font=\footnotesize] at (axis cs:2008,18)
    {物理地板 \~\$20--30/kg};

% Wright's Law 趋势线（示意）
\addplot[color=gray, dotted, line width=0.8pt, mark=none]
    coordinates {
        (1981,54500)(2018,1400)(2026,200)
    };
\node[gray, font=\footnotesize] at (axis cs:1996,12000)
    {Wright 学习率 $\approx$\,85\%};

% 标注
\node[primary, font=\footnotesize] at (axis cs:1996,65000)
    {\$54,500/kg};
\node[primary, font=\footnotesize] at (axis cs:2008,15000)
    {\$13,000/kg};
\node[primary, font=\footnotesize] at (axis cs:2017,3500)
    {\$2,720/kg};
\node[primary, font=\footnotesize] at (axis cs:2021,1100)
    {\$1,400/kg};
\node[accent, font=\footnotesize] at (axis cs:2028,250)
    {\~\$200/kg};

\end{semilogyaxis}
\end{tikzpicture}

\caption{LEO发射成本历史曲线（半对数坐标）：从航天飞机到Starship}
\label{fig:launch_cost}
\end{figure}

从Falcon 9到Starship，预期的\$2,720$\rightarrow$\$200意味着再降93\%——即使假设Wright学习率（每翻倍产量成本降15\%--20\%，与半导体制造经验一致\cite{guangfa2026learning}），达成这一目标也需要Starship的年发射频次从当前每年几次增至数百次。SpaceX的2026年FCC频率申请和S1测试进度\cite{spacex2026fcc,spacex2026s1}表明，高频次复用仍处于早期验证阶段。

\subsection{马斯克"300--500\,GW/年"主张的拆解}

马斯克提出的"Starship每年300\,GW，或许500\,GW"\cite{musk2025starship}值得逐层拆解。

假设条件：Starship单次LEO运力200\,t（FCC申报数据\cite{spacex2026fcc}），每t算力对应约0.5\,MW（基于H100级GPU的功率密度）。则单次发射可部署$\sim$0.1\,GW算力。300\,GW/年意味着每年3,000次Starship发射——即每天8次以上。FAA 2025年对Boca Chica的环评审批\cite{faa2025boca}仅支持年发射25--50次。即便KSC增设发射台\cite{faa2026ksc}，年3000次频次在当前监管、基础设施和制造能力框架下不具备可行性。500\,GW对应的5,000次/年更是脱离了讨论的现实基础。

进一步地：以当前最先进的H100 GPU（700\,W）为例，1\,GW算力对应约143万颗GPU。即使假设Blackwell（1200\,W）\cite{nvidia2026rubin}，也需83万颗。这是一个超出NVIDIA当前全年H100出货量的数量级。算力供给约束——而非仅发射约束——可能才是马斯克主张中最大的隐性瓶颈。

\subsection{轨道太阳能的效率链}

轨道太阳能的计算链需要逐项乘以效率系数：

\[
P_{\text{可用}} = G_{\text{SC}} \times \eta_{\text{PV}} \times \eta_{\text{deg}} \times \eta_{\text{MPPT}} \times \eta_{\text{bat}} \times D_{\text{duty}}
\]

其中：$G_{\text{SC}}=1361$\,\unitWpm{}（太阳常数），$\eta_{\text{PV}}\approx 30\%$（空间级三结砷化镓\cite{zhang2025solar}），$\eta_{\text{deg}}\approx 85\%$（5年退化后），$\eta_{\text{MPPT}}\approx 95\%$，$\eta_{\text{bat}}\approx 88\%$（充放电循环），$D_{\text{duty}}\approx 60\%$（轨道日照占空比）。逐项相乘后，1361\,\unitWpm{}$\rightarrow$约183\,\unitWpm{}可用电力，综合效率仅约13.5\%。

图~\ref{fig:pv_chain} 以瀑布图展示这一逐级退化过程。

\begin{figure}[H]
\centering
% pv_efficiency_chain.tex — 轨道光伏效率链瀑布图
\begin{tikzpicture}[
    node distance=0.6cm,
    block/.style={
        rectangle,
        draw=primary,
        fill=primary!8,
        minimum width=3.0cm,
        minimum height=1.3cm,
        align=center,
        font=\small
    },
    loss/.style={
        rectangle,
        draw=highlight,
        fill=highlight!12,
        minimum width=2.2cm,
        minimum height=0.9cm,
        align=center,
        font=\footnotesize
    },
    arrow/.style={
        ->,
        >=stealth,
        line width=1.0pt,
        color=primary
    },
    label/.style={
        font=\footnotesize,
        text=primary
    }
]

% 主流程：左列
\node[block] (gsc) at (0,0)
    {太阳常数\\[-2pt] {\large\textbf{1,361 W/m\textsuperscript{2}}}};

\node[block] (pv) at (0,-2.8)
    {光伏转换\\[-2pt] {\large\textbf{408 W/m\textsuperscript{2}}}\\（$\eta$=30\% GaAs）};

\node[block] (deg) at (0,-5.6)
    {5年退化后\\[-2pt] {\large\textbf{347 W/m\textsuperscript{2}}}\\（$\eta_{\text{deg}}$=85\%）};

\node[block] (mppt) at (0,-8.4)
    {MPPT+电源管理\\[-2pt] {\large\textbf{330 W/m\textsuperscript{2}}}\\（$\eta_{\text{MPPT}}$=95\%）};

\node[block] (bat) at (0,-11.2)
    {储能充放电\\[-2pt] {\large\textbf{290 W/m\textsuperscript{2}}}\\（$\eta_{\text{bat}}$=88\%）};

\node[block] (duty) at (0,-14.0)
    {日照占空比\\[-2pt] {\large\textbf{174 W/m\textsuperscript{2}}}\\（$D$=60\% LEO）};

\node[block, fill=highlight!20, draw=highlight, line width=1.5pt] (final)
    at (0,-17.0)
    {\textbf{可用电力}\\[-2pt] {\LARGE\textbf{183 W/m\textsuperscript{2}}}\\[-2pt] \footnotesize 综合效率 $\approx$\,13.5\%};

% 损失标注：右列
\node[loss] (loss1) at (3.8,-1.4)
    {损失\\953 W/m\textsuperscript{2}};

\node[loss] (loss2) at (3.8,-4.2)
    {退化损失\\61 W/m\textsuperscript{2}};

\node[loss] (loss3) at (3.8,-7.0)
    {电源损失\\17 W/m\textsuperscript{2}};

\node[loss] (loss4) at (3.8,-9.8)
    {充放损失\\40 W/m\textsuperscript{2}};

\node[loss] (loss5) at (3.8,-12.6)
    {地影损失\\116 W/m\textsuperscript{2}};

% 箭头
\draw[arrow] (gsc) -- (pv);
\draw[arrow] (pv) -- (deg);
\draw[arrow] (deg) -- (mppt);
\draw[arrow] (mppt) -- (bat);
\draw[arrow] (bat) -- (duty);
\draw[arrow] (duty) -- (final);

% 损失箭头
\draw[->, >=stealth, color=highlight, line width=0.8pt]
    (pv.east) -- ++(1.5,0) |- (loss1.west);
\draw[->, >=stealth, color=highlight, line width=0.8pt]
    (deg.east) -- ++(1.5,0) |- (loss2.west);
\draw[->, >=stealth, color=highlight, line width=0.8pt]
    (mppt.east) -- ++(1.5,0) |- (loss3.west);
\draw[->, >=stealth, color=highlight, line width=0.8pt]
    (bat.east) -- ++(1.5,0) |- (loss4.west);
\draw[->, >=stealth, color=highlight, line width=0.8pt]
    (duty.east) -- ++(1.5,0) |- (loss5.west);

% 对比标注
\node[font=\footnotesize, text=gray] at (0,-18.2)
    {对比：地面光伏综合可用系数 $\approx$\,21.3\%};

\end{tikzpicture}

\caption{轨道光伏效率链瀑布图：1361\,W/m\textsuperscript{2} $\rightarrow$ 183\,W/m\textsuperscript{2}可用电力}
\label{fig:pv_chain}
\end{figure}

\begin{table}[H]
\centering
\footnotesize
\caption{轨道与地面光伏效率链对比}
\label{tab:pv_chain}
\begin{tabular}{lcc}
\toprule
\textbf{环节} & \textbf{轨道（三结GaAs）} & \textbf{地面（单晶硅）} \\
\midrule
初始辐照度                    & 1361\,\unitWpm{}         & $\sim$1000\,\unitWpm{} \\
光伏效率                      & $\sim$30\%               & $\sim$24\%              \\
退化后效率（5年）             & $\sim$25.5\%             & $\sim$22\%              \\
MPPT+电源管理                & $\sim$95\%               & $\sim$97\%              \\
储能系统效率                  & $\sim$88\%               & 无（电网）              \\
日照占空比                    & $\sim$60\%               & 无（电网）              \\
\addlinespace
\textbf{综合可用系数}               & $\sim$\textbf{13.5\%}          & $\sim$\textbf{21.3\%}         \\
\bottomrule
\end{tabular}
\end{table}

表~\ref{tab:pv_chain} 显示，尽管轨道太阳常数比地面高36\%，但储能与占空比损耗使得轨道综合可用系数反低于地面光伏。这意味着，马斯克描绘的"无限太阳能"在工程实现中需要搭配显著的光伏阵列面积扩大——进一步加剧上一章讨论的散热面积约束。

\subsection{GPU在轨可靠性：最大的未知}

地面数据中心的GPU年均失效率（AFR）通常在1\%--3\%\cite{uptime2024}。太空高能粒子辐射环境（银河宇宙射线GCR、太阳高能粒子SPE、范艾伦辐射带穿越）会显著增加单粒子效应（SEE）概率\cite{nasa2021tm4527}。NASA对ISS电子设备的跟踪数据\cite{nasa2024iss}及地面辐射模拟\cite{nasajpl2025rad}表明，无抗辐射加固的商业级GPU在LEO轨道的AFR预计为8\%--15\%，极端太阳事件下甚至可达30\%+。

图~\ref{fig:gpu_afr} 将这一差异以柱状图可视化。

\begin{figure}[H]
\centering
% gpu_reliability_bars.tex — GPU 可靠性对比柱状图
\begin{tikzpicture}
\begin{axis}[
    width=0.95\textwidth,
    height=6.5cm,
    ybar,
    bar width=28pt,
    ylabel={年化失效率 AFR（\%）},
    ymin=0, ymax=38,
    ytick={0,3,5,10,15,30},
    yticklabels={0,3,5,10,15,30},
    symbolic x coords={地面数据中心, 轨道商业级GPU, 轨道抗辐射GPU, 极端太阳事件},
    xtick=data,
    xticklabel style={align=center, font=\footnotesize},
    enlarge x limits=0.25,
    nodes near coords,
    nodes near coords style={font=\footnotesize},
    point meta=explicit symbolic,
    every node near coord/.append style={yshift=6pt},
    legend style={
        at={(0.5,-0.22)},
        anchor=north,
        legend columns=1,
        draw=none,
        font=\footnotesize
    },
    grid=major,
    grid style={draw=gray!20},
    axis lines=left,
    enlarge y limits={upper, value=0.2}
]

% 柱状图数据
\addplot[fill=primary!60, draw=primary]
    coordinates {
        (地面数据中心, 2) [2\%]
        (轨道商业级GPU, 12) [12\%]
        (轨道抗辐射GPU, 5) [5\%]
        (极端太阳事件, 30) [30\%+]
    };

% 标注
\draw[->, >=stealth, highlight, line width=1.0pt, dashed]
    (axis cs:地面数据中心,8) -- (axis cs:轨道商业级GPU,8)
    node[midway, above, font=\footnotesize, text=highlight] {$\times$6};

\draw[->, >=stealth, highlight, line width=1.0pt, dashed]
    (axis cs:轨道商业级GPU,20) -- (axis cs:极端太阳事件,20)
    node[midway, above, font=\footnotesize, text=highlight] {$\times$2.5};

% 不可接受阈值线
\draw[red!50, dashed, line width=0.8pt]
    (axis cs:地面数据中心,15) -- (axis cs:极端太阳事件,15);
\node[red!60, font=\footnotesize] at (axis cs:轨道抗辐射GPU,16.5)
    {不可接受阈值};

% 注解
\node[align=left, font=\footnotesize, text=gray]
    at (axis cs:地面数据中心,36)
    {在轨维修 $\rightarrow$ 不可行\\集群3年可用算力 $\rightarrow$ 20--40\%};

\end{axis}
\end{tikzpicture}

\caption{GPU年化失效率对比：地面 vs 轨道（LEO）}
\label{fig:gpu_afr}
\end{figure}

\begin{table}[H]
\centering
\footnotesize
\caption{GPU可靠性对比：轨道 vs 地面}
\label{tab:gpu_reliability}
\begin{tabular}{lcc}
\toprule
\textbf{指标} & \textbf{轨道（LEO）} & \textbf{地面数据中心} \\
\midrule
主要辐射源 & GCR, SPE, 范艾伦带 & 本底辐射（可忽略） \\
GPU AFR（商业级，无加固） & 8\%--15\%（估计） & 1\%--3\%\cite{uptime2024} \\
极端事件下AFR & 30\%+（一级耀斑） & 不可比 \\
在轨更换/维修 & \textbf{不可行}（当前无在轨服务能力） & 热插拔，<1小时 \\
抗辐射加固代价 & 3--5倍成本 + 30--50\%性能损失\cite{gaoanxin2025rad} & 不适用 \\
\bottomrule
\end{tabular}
\end{table}

表~\ref{tab:gpu_reliability} 揭示了轨道算力最关键的单点故障：**在轨维护的不可行性**。即使GPU年失效率仅提高至地面2--3倍（保守假设5\%--8\%），在无法更换的条件下，集群可用算力将在3年内衰减至初始的20\%--40\%。这与地面数据中心"坏一块换一块"的运维模式构成根本性差异。NASA的Katalyst计划\cite{nasa2026katalyst}正在开发在轨服务机器人，但预计最早在2027--2028年才有初步演示能力。黄仁勋在Q4财报中强调的"将需要数年"\cite{huang2026earnings}，很大程度上正是指向这些非芯片层面的系统工程缺口。

\subsection{工程层小结}

发射成本是四大参数中最乐观的一项——Starship确有降本潜力，但年3000次频次缺乏现实支撑。轨道光伏综合效率反低于地面，供电优势不像表面数字那么直观。最关键的是：GPU在轨可靠性数据几近空白，缺乏在轨维护能力的条件下，集群寿命约束将使TCO中的硬件替换成本急剧膨胀。黄仁勋的"硬件工程谨慎"在这一层获得最充分的证据支持。
